\section{Path-oracle equivalence under sample splitting}
\label{sec:supp-oracle}

Let \(\mathcal T_k\) collect the fitting-fold and selection-fold data associated with evaluation fold \(k\).
Write \(\widehat F_{h,k}(b)\) for the evaluation-fold estimator formed with augmentation path \(b\).

\begin{lemma}[Regression-path perturbation]
\label{lem:nuisance-perturb}
Under Conditions~\ref{cond:sampling}--\ref{cond:entropy},
\[
\sqrt{N_{M,h}}
\left\|
\widehat F_{h,k}(a_k^{\mathrm{or}})
-
\widehat F_{h,k}(a_M)
\right\|_{\mathcal Y}
=o_{\mathbb P}(1),
\]
provided
\[
\|a_k^{\mathrm{or}}-a_M\|_{\mathrm{ver}}
=o_{\mathbb P}(1)
\]
and the realized difference class satisfies Condition~\ref{cond:entropy}.
\end{lemma}

\begin{proof}
For every \(y\),
\begin{equation}
U_y(a_k^{\mathrm{or}})-U_y(a_M)
=
\{a_{k,y}^{\mathrm{or}}-a_{M,y}\}
\left(1-\frac{R}{\pi_M}\right).
\label{eq:nuisance-exact}
\end{equation}
Conditional on \(\mathcal T_k\), the right side has conditional mean zero given \(W\). Define
\[
\mathbb D_{M,k,y}
=
\sqrt{p_Mh^d}
\frac{K_h}{\mu_h}
\{a_{k,y}^{\mathrm{or}}-a_{M,y}\}
\left(1-\frac{R}{\pi_M}\right).
\]
Its conditional squared \(L_2(P)\) norm is
\[
\frac{p_Mh^d}{\mu_h^2}
E\left[
K_h^2
\left(\frac1{\pi_M}-1\right)
\{a_{k,y}^{\mathrm{or}}-a_{M,y}\}^2
\,\Bigm|\,\mathcal T_k
\right],
\]
which is uniformly \(o_{\mathbb P}(1)\). The conditional maximal inequality for VC-type classes gives
\[
\left\|
\frac1{\sqrt{|I_{E,k}|}}
\sum_{i\in I_{E,k}}
\mathbb D_{M,k,\cdot}(O_i)
\right\|_{\mathcal Y}
=o_{\mathbb P}(1).
\]
Because \(|I_{E,k}|/M\to\rho_k\in(0,1)\), replacing \(\sqrt{|I_{E,k}|}\) by \(\sqrt M\) changes only a bounded deterministic factor. Lemma~\ref{lem:ratio} then gives the stated estimator difference.
\end{proof}

\begin{lemma}[Coefficient perturbation]
\label{lem:coefficient-perturb}
Under Conditions~\ref{cond:sampling}--\ref{cond:path-compatibility},
\[
\sqrt{N_{M,h}}
\left\|
\widehat F_{h,k}(\widehat a_k)
-
\widehat F_{h,k}(a_k^{\mathrm{or}})
\right\|_{\mathcal Y}
=o_{\mathbb P}(1),
\]
where
\[
a_{k,y}^{\mathrm{or}}
=
\widehat q_{k,y}
+
\lambda_{h,k}^{\mathrm{or}}
\widehat D_{k,y},
\qquad
\widehat D_{k,y}=\widehat m_{k,y}-\widehat q_{k,y},
\]
uses the same fitted regressions as \(\widehat a_k\) and replaces only the estimated coefficient by its conditional population minimizer.
\end{lemma}

\begin{proof}
Conditional on \(\mathcal T_k\),
\begin{equation}
U_y(\widehat a_k)-U_y(a_k^{\mathrm{or}})
=
(\widehat\lambda_{h,k}-\lambda_{h,k}^{\mathrm{or}})
\widehat D_{k,y}(W)
\left(1-\frac{R}{\pi_M(W)}\right).
\label{eq:lambda-exact}
\end{equation}
The indexed factor following the coefficient difference is a centered local empirical process. Let \(\widehat{\mathcal D}_{h,k}^2\) be its fitted squared maximal local \(L_2\) radius. Condition~\ref{cond:path-compatibility} gives
\[
\widehat{\mathcal D}_{h,k}^2
\le
\omega(\mathcal H_{h,k})+o_{\mathbb P}(1).
\]
The conditional maximal inequality and Lemma~\ref{lem:ratio} yield
\[
\sqrt{N_{M,h}}
\left\|
\widehat F_{h,k}(\widehat a_k)
-
\widehat F_{h,k}(a_k^{\mathrm{or}})
\right\|_{\mathcal Y}
=
O_{\mathbb P}\left(
|\widehat\lambda_{h,k}-\lambda_{h,k}^{\mathrm{or}}|
\widehat{\mathcal D}_{h,k}
\right)+o_{\mathbb P}(1).
\]
If \(\mathcal H_{h,k}\to H_{\infty,k}>0\), Lemma~\ref{lem:regret} gives
\[
|\widehat\lambda_{h,k}-\lambda_{h,k}^{\mathrm{or}}|
=o_{\mathbb P}(1),
\]
and \(\widehat{\mathcal D}_{h,k}=O_{\mathbb P}(1)\). If \(\mathcal H_{h,k}\to0\), Condition~\ref{cond:path-compatibility} gives \(\widehat{\mathcal D}_{h,k}\to_{\mathbb P}0\), while the coefficient difference is bounded by one. The product is \(o_{\mathbb P}(1)\) in both regimes.
\end{proof}

\begin{proof}[Proof of Theorem~\ref{thm:oracle-equivalence}]
The gain-estimation rates and criterion results follow from Lemmas~\ref{lem:gain-rates} and \ref{lem:regret}. Lemma~\ref{lem:coefficient-perturb} replaces the estimated coefficient by the fold-specific conditional population minimizer while holding the fitted regressions fixed. Weighted summation over the three evaluation folds preserves the \(o_{\mathbb P}(1)\) order. Lemma~\ref{lem:nuisance-perturb} then replaces the fitted population-minimizing paths by the common triangular path \(a_M\). The positive-curvature coefficient rate and the interior quadratic regret are the corresponding statements of Lemma~\ref{lem:regret}; the vanishing-curvature result follows from Condition~\ref{cond:path-compatibility}.
\end{proof}
